\documentclass[10pt]{letter}
\usepackage[a4paper, margin=0.85in]{geometry}
\usepackage[T1]{fontenc}
\setlength{\parskip}{5pt}
\signature{Harsh Bordekar \\ harshbordekar@gmail.com}
\begin{document}
\begin{letter}{Helsing GmbH \\ Munich, Germany}
\opening{Dear Hiring Team,}

I am writing to apply for your Simulation Engineer position on the Physical Products team. I am a simulation and structural analysis engineer with over seven years building, verifying, and maintaining physics-based simulation environments for real aerospace structures and vehicle hardware, work that has always had to earn its way onto a test article rather than replace one.

At DLR and TU Braunschweig I built and maintain a probabilistic, multiscale simulation environment for leakage onset in CFRP hydrogen storage tanks, coupling microstructure-informed damage models with stochastic defect distributions to propagate failure from ply to structural scale. I developed Monte Carlo uncertainty quantification and sensitivity-analysis capability directly into that work, generating probability-of-failure relationships across the real range of strain and load conditions rather than a single nominal case, and I derived the certification-aligned failure criterion that turns those simulation outputs into structural sizing inputs for design trade studies. I also ran large-scale thermal and structural simulation campaigns on HPC/SLURM infrastructure for the Callisto reusable-rocket programme's metal upper housing bay.

Model credibility is something I treat as a deliverable, not an afterthought. I verify and technically review externally and student-produced simulation models against defined acceptance criteria, tracing failure modes to numerical, meshing, or data-quality issues, and I have built and validated models against real measured and test data throughout my career: fatigue and defect-tolerance work correlated against measured defect populations, and a CT-based defect-detection classifier validated to 91\% accuracy against ground-truth labels. I also build tooling that makes simulation faster and more accessible to a wider team, including an agentic surrogate application that replaces full-scale FEM homogenization runs with on-demand inference. Beyond model-building, I write performance-sensitive C++ for FEM post-processing and work daily in Python, MATLAB, and ABAQUS/Nastran.

I also teach: I instruct 60+ Master's students each year in FEM theory and ABAQUS-based composite modelling, which has sharpened my ability to explain simulation results and their limitations clearly to technical and non-technical audiences alike.

To be transparent about fit: my simulation background is in structures, thermal analysis, and damage/fracture mechanics rather than flight dynamics, aerodynamics, or GNC, and I have not worked on 6-DOF vehicle or engagement-level simulation, CFD, or real-time/HITL environments. My validation experience is at the component and lab-test level (defect measurements, CT ground truth, external model review) rather than flight or range test data. What I bring is a genuine, hands-on habit of building simulation environments that people other than me trust enough to make a design decision on, and the physics-first discipline to know when a model is a reasonable simplification and when it isn't.

I am motivated by Helsing's mission and by the scale of the engineering problem behind it, and I would welcome the chance to bring my simulation and model-validation background into your Physical Products team.

Thank you for your consideration. I look forward to discussing how I can contribute.

\closing{Sincerely,}
\end{letter}
\end{document}
